\hypertarget{repair-rounds-scheduler-variants-and-the-axial-handoff}{%
\chapter{Repair rounds, scheduler variants and the axial
handoff}\label{repair-repair-rounds-scheduler-variants-and-the-axial-handoff}}

\begin{quote}
Follow-up (\texttt{docs/megascale.md}): scaling to 1e4, 1e5 and 1e6
sites. A 32 x 4 engine block assembles 1e6 sites in 1.5 s and holds them
at the loss-limited floor rho \textasciitilde{} 1.5-1.8 (eps\_img +
T\_img/tau); static trapping power, the objective\textquotesingle s
field and the 3 um lane distance are the walls.
\end{quote}

Follow-up to \texttt{docs/rearrangement.md}. Code:
\texttt{picsim/rearrange.py} (\texttt{repair},
\texttt{cycle\_with\_repair}, selection modes),
\texttt{scripts/m21\_repair.py},
\texttt{scripts/m19\_quantum.py\ axial}; results
\texttt{results/m21\_*.json}, \texttt{results/m19\_axial.json}. Values
are SIM and valid in the regime stated. They should be reviewed before
they are quoted.

\hypertarget{1-result}{%
\section{Result}\label{repair-1-result}}

After the rearrangement cycle (32 x 32 block from a 50 x 45 array at
50\% loading, 17 ms), one round of imaging and direct repair takes 10 to
13 ms: a 10 ms image and 0.1 to 1.3 repair passes. Each repair moves a
single reservoir atom through the lanes into a hole. The repair works,
but it cannot get below a floor set by the image itself. Atoms lost
because of the image are not seen by it, and every atom is also exposed
to vacuum loss for the round\textquotesingle s duration. The floor is

\begin{verbatim}
defects ~ N (eps_img + t_round / tau) + p_move x (repairs),   N = 1024,
\end{verbatim}

about 0.14 defects for eps\_img = 1e-4 and tau = 300 s. A second round
does not lower it. Hence the probability of a defect-free 32 x 32 block
is set by the imaging loss and the vacuum lifetime, not by the
transport, once the per-move loss is about 1e-4 or less:

{\small
\begin{longtable}[]{@{}>{\raggedright\arraybackslash}p{0.112\dimexpr(\linewidth-10\tabcolsep)\relax}>{\raggedright\arraybackslash}p{0.129\dimexpr(\linewidth-10\tabcolsep)\relax}>{\raggedright\arraybackslash}p{0.103\dimexpr(\linewidth-10\tabcolsep)\relax}>{\raggedright\arraybackslash}p{0.466\dimexpr(\linewidth-10\tabcolsep)\relax}>{\raggedright\arraybackslash}p{0.190\dimexpr(\linewidth-10\tabcolsep)\relax}@{}}
\toprule\noalign{}
per-move loss & vacuum lifetime & imaging loss & P(defect-free): after
cycle / after 1 repair / after 2 & defects after 1 repair \\
\midrule\noalign{}
\endhead
\bottomrule\noalign{}
\endlastfoot
1e-3 & 300 s & 1e-4 & 0.27 / \textbf{0.87} / 0.87 & 0.14 \\
1e-4 & 300 s & 1e-4 & \textbf{0.84} / 0.87 / 0.87 & 0.14 \\
1e-3 & 300 s & 1e-3 & 0.27 / 0.35 / 0.33 & 1.04 \\
1e-4 & 300 s & 1e-3 & \textbf{0.84 / 0.35} / 0.33 & 1.03 \\
1e-3 & 30 s & 1e-4 & 0.14 / 0.59 / 0.62 & 0.53 \\
1e-4 & 30 s & 1e-4 & 0.48 / 0.62 / 0.63 & 0.48 \\
\end{longtable}
}

(40 loadings x 25 loss draws per row.) With good transport (1e-4) and an
imaging loss of 1e-3, a repair round makes things worse: 0.84 becomes
0.35, because the image costs about one atom of the block. Imaging
should then be skipped, or used only to post-select. For a 99\%
defect-free block one needs roughly 1e-5 for each term: per-move loss,
imaging loss per image, and t\_round / tau. These are properties of the
platform, not of the photonic chip. The chip\textquotesingle s
contribution is 1.1 moves per atom and 17 ms per cycle.

\hypertarget{2-scheduler-variants}{%
\section{Scheduler variants}\label{repair-2-scheduler-variants}}

Balancing-pass selection on 40 loadings (32 x 32, 50\%, one stack):

{\small
\begin{longtable}[]{@{}>{\raggedright\arraybackslash}p{0.638\dimexpr(\linewidth-6\tabcolsep)\relax}>{\raggedright\arraybackslash}p{0.128\dimexpr(\linewidth-6\tabcolsep)\relax}>{\raggedright\arraybackslash}p{0.234\dimexpr(\linewidth-6\tabcolsep)\relax}@{}}
\toprule\noalign{}
mode & passes (all) & cycle time, mean / max \\
\midrule\noalign{}
\endhead
\bottomrule\noalign{}
\endlastfoot
greedy (most atoms per pass) & 8.05 & 17.8 / 23.1 ms \\
rate (most atoms per unit pass time) & 8.05 & 17.6 / 22.8 ms \\
lookahead (best two-pass sequence among the 8 best first passes) & 7.98
& 17.6 / 22.9 ms \\
\end{longtable}
}

The three modes agree to within 1\%. The balancing stage has a floor of
about four passes, set by the geometry. In one pass the 16-row window
can hold at most 9 non-target rows above or below the 32-row target
band, so it can feed at most 9 target rows. Feeding all 32 therefore
takes at least four passes, plus two compaction passes. The greedy
schedule (5.8 + 2) is within about 30\% of that floor. Further gains
would come from shorter passes, not fewer: ramps, lane steps and scanner
settling make up 1.3 of the 2.2 ms per pass.

\hypertarget{3-axial-motion-during-the-handoff}{%
\section{Axial motion during the
handoff}\label{repair-3-axial-motion-during-the-handoff}}

1D split-step propagation along the optical axis from the ground state
of the static site (axial trap quantum 0.40 uK at 0.1 mK). The static
site\textquotesingle s depth follows the lateral separation of the
moving-trap centre, which carries the atom 2 um into the lane and back,
and the moving trap\textquotesingle s axial profile comes from the 3D
grid. The parameter is the axial offset dz between the 1075 nm and 1064
nm foci (results/m19\_axial.json):

{\footnotesize
\begin{longtable}[]{@{}llllll@{}}
\toprule\noalign{}
hold / moving & ramp / lane step & dz = 0 & 0.5 um & 1 um & 2 um \\
\midrule\noalign{}
\endhead
\bottomrule\noalign{}
\endlastfoot
0.1 / 0.3 mK & 100 / 100 us & 0.06 uK & 4.1 uK & 16.6 uK & 40.7 uK \\
0.1 / 0.3 mK & 200 / 400 us & \textless{} 0.01 uK & \textbf{0.19 uK} &
2.0 uK & 8.3 uK \\
0.1 / 0.5 mK & 100 / 100 us & 0.02 uK & 6.4 uK & 31.0 uK & 28.3 uK \\
0.1 / 0.5 mK & 200 / 400 us & \textless{} 0.01 uK & \textbf{0.22 uK} &
1.8 uK & 4.4 uK \\
\end{longtable}
}

(Quantum energy gain; the classical Wigner values agree to within 20\%
above 1 uK and are pessimistic below.) With the foci coincident, the
axial motion is not excited: the change of axial trap frequency during
the handoff is adiabatic. A focus offset displaces the axial equilibrium
twice per move, and that is what heats. With the slow handoff, keeping
the heating per move below the lateral 0.12 to 0.15 uK of
\texttt{docs/rearrangement.md} needs dz \textless= 0.5 um. A 1 um offset
costs about 2 uK per move, and 2 um costs 4 to 8 uK. This tightens the
+/-1 um delivery tolerance of \texttt{docs/transport3d.md} into a
heating tolerance of +/-0.5 um. Since a diffractive objective puts the
two foci about 170 um apart, the 1064 nm path needs a focus adjustment
that holds 0.5 um (0.3\% of the chromatic offset), for example a slow
stage on the relay\textquotesingle s last y and x lenses, calibrated on
the atoms.

\hypertarget{4-limits-and-next-steps}{%
\section{Limits and next steps}\label{repair-4-limits-and-next-steps}}

The repair model moves one reservoir atom per hole, through the lanes;
repaired atoms come from the nearest rows under the pass rules.
Detection is assumed perfect apart from the imaging-induced loss, with
no false positives or negatives. The imaging time (10 ms) and losses are
CHOICE values, not measured. The axial check is 1D, with the atom
assumed to sit at the moving-trap centre laterally.

Next steps: (i) image-free or low-loss detection (the floor of Sec. 1);
(ii) shorter passes (Sec. 2); (iii) an active focus lock between the
1064 and 1075 nm paths to 0.5 um (Sec. 3).

\hypertarget{verification-notes}{%
\section{Verification notes}\label{repair-verification-notes}}

{\small
\begin{longtable}[]{@{}>{\raggedright\arraybackslash}p{0.504\dimexpr(\linewidth-8\tabcolsep)\relax}>{\raggedright\arraybackslash}p{0.101\dimexpr(\linewidth-8\tabcolsep)\relax}>{\raggedright\arraybackslash}p{0.076\dimexpr(\linewidth-8\tabcolsep)\relax}>{\raggedright\arraybackslash}p{0.319\dimexpr(\linewidth-8\tabcolsep)\relax}@{}}
\toprule\noalign{}
value & location & type & pointer \\
\midrule\noalign{}
\endhead
\bottomrule\noalign{}
\endlastfoot
repair table & Sec. 1 & SIM & results/m21\_repair.json (review) \\
scheduler modes & Sec. 2 & SIM & results/m21\_modes.json \\
axial quantum and Wigner results & Sec. 3 & SIM &
results/m19\_axial.json (review) \\
eps\_img 1e-3 / 1e-4, t\_img 10 ms, tau 30 / 300 s, p\_move 1e-3 / 1e-4
& Sec. 1 & CHOICE & assumptions \\
\end{longtable}
}
